%! Multiplying \& Dividing Radicals; Rationalizing — Unit 5, Lesson 5.3 — Warm-Up (Answer Key)
\documentclass[10pt]{article}
\usepackage{algebra2-article}
\usepackage{algebra2-key}   % pulls in -boxes; do NOT also load -boxes
\newcommand{\TallMath}[1]{$\displaystyle #1\rule[-1.4em]{0pt}{3.2em}$}
% In-formula answer slot (\ans is text-mode and must never appear inside $...$).
\newcommand{\mans}[1]{{\color{keyred}\mathbf{#1}}}

\begin{document}

\pageheader{Unit 5, Lesson 5.3}{Warm-Up --- Answer Key}
\namedateperiod

\vspace{0.10in}

\begin{notesbox}{Getting Ready: Skills We'll Use Today}
\small Three skills you already have. Today they turn into one new one. Work quickly.

\vspace{4pt}
\textbf{1. Distributing and multiplying binomials.} Expand and simplify.
\begin{center}
\renewcommand{\arraystretch}{1.6}
\begin{tabular}{@{}l l l@{}}
(a) \; $3x(2x+5)=$ \ans{$6x^{2}+15x$} & (b) \; $(x+4)(x-4)=$ \ans{$x^{2}-16$} &
(c) \; $(x+4)^{2}=$ \ans{$x^{2}+8x+16$} \\
\end{tabular}
\end{center}
In (b) two of the four terms cancelled and the answer had no middle term. What was special about
those two binomials? \ansline{They had the \emph{same two terms} with the \emph{opposite} middle
sign --- $+4$ in one, $-4$ in the other. The outer and inner products came out as $-4x$ and $+4x$,
so they cancelled, leaving only $x^{2}-16$.}

\vspace{4pt}
\textbf{2. Yesterday's radicals.} Simplify each.
\begin{center}
\renewcommand{\arraystretch}{1.6}
\begin{tabular}{@{}l l l l@{}}
(a) \; $\sqrt{12}=$ \ans{$2\sqrt3$} & (b) \; $\sqrt{50}=$ \ans{$5\sqrt2$} &
(c) \; $\sqrt{7}\cdot\sqrt{7}=$ \ans{$7$} & (d) \; $\left(\sqrt{11}\right)^{2}=$ \ans{$11$} \\
\end{tabular}
\end{center}
Look at (c) and (d). What does multiplying a square root \emph{by itself} do to the radical sign?
\ansline{It erases it. $\sqrt7\cdot\sqrt7=\sqrt{49}=7$ --- squaring is exactly the operation a
square root undoes, so the radical disappears and a plain rational number is left behind.}

\vspace{4pt}
\textbf{3. Multiplying by a form of $1$.}
\begin{itemize}[leftmargin=1.5em, itemsep=3pt]
  \item Fill in the missing numerator: \; $\dfrac{3}{4}=\dfrac{\mans{15}}{20}$. \;
        What did you multiply the top \emph{and} the bottom by? \ans{$5$}
  \item Multiplying a fraction by $\dfrac{5}{5}$ does not change its value, because
        $\dfrac{5}{5}=$ \ans{$1$}. \; Is the same true of $\dfrac{\sqrt{2}}{\sqrt{2}}$?
        \ans{yes --- it is also $1$}
\end{itemize}

\end{notesbox}

\begin{teachernote}
Four minutes. Each item is one third of today's lesson, disassembled.

\textbf{Item 1} is the entire multiplication half. Radicals distribute and FOIL exactly the way
these binomials do --- $\sqrt3\left(\sqrt6+\sqrt{15}\right)$ is $3x(2x+5)$ in different clothing.
Item 1(b) is the one to debrief carefully: get the words \emph{``same terms, opposite middle
sign''} said aloud, because that sentence is the definition of a \textbf{conjugate}, which arrives
in Section 3 and runs the second half of the period. Item 1(c) plants the middle-term error before
it can happen; when a student later writes $\left(\sqrt5+\sqrt2\right)^{2}=7$, point back at
$(x+4)^{2}\ne x^{2}+16$, which they have already rejected in this room.

\textbf{Item 2(c)--(d) is the engine of the whole lesson.} $\sqrt a\cdot\sqrt a=a$ is the only new
mechanical idea today: it is what makes a product come out rational, and it is why multiplying a
denominator by its own radical clears it. Make sure the word \emph{erases} (or ``undoes,'' or
``cancels the radical'') gets said out loud now.

\textbf{Item 3} is the honesty check on rationalizing. Students often suspect that multiplying by
$\frac{\sqrt2}{\sqrt2}$ is cheating --- that it changes the number. It does not, for exactly the
reason $\frac{3}{4}=\frac{15}{20}$: multiplying by a fraction whose top and bottom match is
multiplying by $1$. The form changes; the value cannot. Anyone who hesitates on the last blank will
be the student who mistrusts Section 5.
\end{teachernote}

\end{document}
